% =====================================================================
%  2bat-ud5-15.tex  —  SESSIÓ 15: Taller de preparació de la prova
%  (sense preàmbul: s'inclou des de main.tex)
% =====================================================================
\horatitol{Sessió 15 — Taller de preparació de la prova}%
{Conèixer l'estructura de la prova i la rúbrica, assajar un simulacre per blocs i fer un pla d'estudi personal.}

\subsection{Idea clau}

Ens preparem per a la prova \textbf{sense sorpreses}. La prova de la sessió següent té
\textbf{quatre blocs}, un per cada criteri (C1--C4). Saber què es demana i com es
corregeix fa que l'examen deixi de ser imprevisible. Després assajarem amb un
\textbf{simulacre} amb temps controlat.

\begin{idea}
\textbf{Estructura de la prova} (un bloc per criteri):
\begin{enumerate}[leftmargin=1.6em, itemsep=2pt]
  \item \textbf{Bloc 1 (C1) — Probabilitat condicionada:} llegir una taula i calcular
        $P(A\,|\,B)$ amb la base correcta.
  \item \textbf{Bloc 2 (C2) — Binomial:} calcular $P(k)$ i el nombre esperat $\mu=n\cdot p$.
  \item \textbf{Bloc 3 (C3) — Normal (tipificar):} passar a $Z$ i llegir la taula.
  \item \textbf{Bloc 4 (C4) — Normal (intervals i decisió):} calcular $P(a<X<b)$ i
        interpretar-ho en context.
\end{enumerate}
\textbf{Es valora}, a més del resultat: triar la \textbf{base correcta}, reconèixer si la
situació és \textbf{binomial}, tipificar bé i \textbf{interpretar} en el context social.
\end{idea}

\subsection{Rúbrica simplificada}

\noindent Així es corregeix cada bloc (nivells d'assoliment):
\begin{center}
\renewcommand{\arraystretch}{1.3}
\begin{tabular}{p{2.4cm} p{10.2cm}}
\toprule
\textbf{Nivell} & \textbf{Què demostra} \\
\midrule
\textbf{AE} (excel·lent) & Resol sense errors i interpreta el resultat en el context
social. \\
\textbf{AN} (notable) & Resol correctament amb algun error menor; la interpretació és
essencialment bona. \\
\textbf{AS} (suficient) & Resol el procediment bàsic amb alguns errors; interpretació
incompleta. \\
\textbf{NA} (no assolit) & No identifica el procediment o comet errors que
n'invaliden el resultat. \\
\bottomrule
\end{tabular}
\end{center}

\subsection{Simulacre — un ítem de cada bloc}

\begin{exemple}[com es resol un ítem ben fet]
\textbf{Ítem tipus (Bloc 3):} en una $N(50,10)$, calcula $P(X<65)$. Un \textbf{ítem
complet} tipifica, llegeix la taula i \textbf{interpreta}:
\[
Z=\frac{65-50}{10}=1,5\ \Rightarrow\ P(X<65)=P(Z<1,5)=0,9332=93,3\%.
\]
\emph{Escriure la tipificació i el valor de la taula, no només el resultat final,
distingeix un AE d'un AS.}
\end{exemple}

\begin{exercicis}
\textbf{Resol el simulacre amb temps controlat. Un ítem per bloc:}
\begin{exlist}
  \item \textbf{(Bloc 1 — C1)} En una taula, $90$ usuaris són joves i, d'aquests, $54$ han
        completat un programa. Calcula $P(\text{completa}\,|\,\text{jove})$ i expressa-la
        en \%.
  \item \textbf{(Bloc 2 — C2)} Cada usuari té un $p=0,6$ de superar una prova. En un grup
        de $10$, calcula el nombre esperat i la probabilitat que en superin
        \textbf{exactament 7}. (Pista: $\binom{10}{7}=120$.)
  \item \textbf{(Bloc 3 — C3)} En una $N(70,8)$, calcula $P(X<78)$ tipificant. (Pista:
        $P(Z<1)=0,8413$.)
  \item \textbf{(Bloc 4 — C4)} El temps d'atenció segueix $N(25,5)$ minuts. Calcula
        $P(20<X<30)$ i interpreta el resultat.
\end{exlist}
\end{exercicis}

\subsection{Formulari-resum (suport)}

\begin{idea}
\textbf{Tingues-ho a mà mentre prepares la prova:}
\begin{itemize}[leftmargin=1.4em, itemsep=2pt]
  \item \textbf{Condicionada:} $P(A\,|\,B)$ divideix pel total \textbf{del grup $B$}.
  \item \textbf{Binomial:} $P(k)=\binom{n}{k}p^{k}(1-p)^{n-k}$; nombre esperat
        $\mu=n\cdot p$.
  \item \textbf{Tipificar:} $Z=\dfrac{x-\mu}{\sigma}$; la taula dona $P(Z<z)$ (a
        l'esquerra).
  \item \textbf{Intervals:} $P(a<X<b)=P(Z<z_b)-P(Z<z_a)$; a la dreta, $1-P(Z<z)$.
  \item \textbf{Regla 68--95--99,7:} $\mu\pm\sigma\to 68\%$; $\mu\pm2\sigma\to 95\%$;
        $\mu\pm3\sigma\to 99,7\%$.
\end{itemize}
\end{idea}

\noindent\textbf{Pla d'estudi personal.} A partir de la graella d'autoavaluació (sessió 3)
i del simulacre, anota els \textbf{dos o tres aspectes} que has de repassar amb més
atenció:
\begin{center}
\renewcommand{\arraystretch}{1.5}
\begin{tabular}{c p{10.5cm}}
\toprule
\textbf{\#} & \textbf{Què he de repassar (i com)} \\
\midrule
1 & \rule{9.8cm}{0.4pt} \\
2 & \rule{9.8cm}{0.4pt} \\
3 & \rule{9.8cm}{0.4pt} \\
\bottomrule
\end{tabular}
\end{center}

% ---------------------------------------------------------------------
\fullsolucions{Sessió 15 — simulacre}
\begin{solucions}
\begin{sollist}
  \item \textbf{(Bloc 1)} Base: joves ($90$); completen $54$.
        $P(\text{completa}\,|\,\text{jove})=\dfrac{54}{90}=\dfrac{3}{5}=0,6=60\%$.
  \item \textbf{(Bloc 2)} $\mu=10\cdot 0,6=6$. $P(7)=\binom{10}{7}(0,6)^7(0,4)^3=
        120\cdot 0,0279936\cdot 0,064\approx 0,215=21,5\%$.
  \item \textbf{(Bloc 3)} $Z=\dfrac{78-70}{8}=1$; $P(X<78)=P(Z<1)=0,8413=84,13\%$.
  \item \textbf{(Bloc 4)} $z_{20}=\dfrac{20-25}{5}=-1$, $z_{30}=\dfrac{30-25}{5}=1$;
        $P(20<X<30)=0,8413-0,1587=0,6826\approx 68,3\%$. Aproximadament el $68\%$ de les
        atencions duren entre $20$ i $30$ minuts (l'interval $\mu\pm\sigma$).
\end{sollist}
\end{solucions}
